\chapter{2022年全国乙卷(理)}

\section{单选题}

\begin{problem}
设全集 \(U=\{1,2,3,4,5\}\),集合 \(M\) 满足 \(\complement_U M=\{1,3\}\),则(\quad)

\choices
    {\(2\in M\)}
    {\(3\in M\)}
    {\(4\notin M\)}
    {\(5\notin M\)}
\end{problem}

\begin{problem}
已知 \(z=1-2\i\),且 \(z+a\overline z+b=0\),其中 \(a\),\(b\) 为实数,则(\quad)

\choices
    {\(a=1\),\(b=-2\)}
    {\(a=-1\),\(b=2\)}
    {\(a=1\),\(b=2\)}
    {\(a=-1\),\(b=-2\)}
\end{problem}

\begin{problem}
已知向量 \(\bs{a}\),\(\bs{b}\) 满足 \(|\bs{a}|=1\),\(|\bs{b}|=\sqrt{3}\),\(|\bs{a}-2\bs{b}|=3\),则 \(\bs{a}\cdot\bs{b}=\)(\quad)

\choices
    {\(-2\)}
    {\(-1\)}
    {\(1\)}
    {\(2\)}
\end{problem}

\begin{problem}
嫦娥二号卫星在完成探月任务后,继续进行深空探测,成为我国第一颗环绕太阳飞行的人造行星,为研究嫦娥二号绕日周期与地球绕日周期的比值,用到数列 \(\{b_n\}\):\(b_1=1+\frac{1}{\alpha_1}\),\(b_2=1+\frac{1}{\alpha_1+\frac{1}{\alpha_2}}\),\(b_3=1+\frac{1}{\alpha_1+\frac{1}{\alpha_2+\frac{1}{\alpha_3}}}\),…,依此类推,其中 \(\alpha_k\in\N^*\)(\(k=1\),\(2\),\(\dots\)),则(\quad)

\choices
    {\(b_1<b_5\)}
    {\(b_3<b_8\)}
    {\(b_6<b_2\)}
    {\(b_4<b_7\)}
\end{problem}

\begin{problem}
设 \(F\) 为抛物线 \(C:y^2=4x\) 的焦点,点 \(A\) 在 \(C\) 上,点 \(B(3,0)\),若 \(AF=BF\),则 \(AB=\)(\quad)

\choices
    {\(2\)}
    {\(2\sqrt{2}\)}
    {\(3\)}
    {\(3\sqrt{2}\)}
\end{problem}

\begin{problem}
执行下边的程序框图,输出的 \(n=\)(\quad)

\bitmapfigure[max width=6.0cm]{img/2022/national_paper_b_science/q6_fig1.png}

\choices
    {\(3\)}
    {\(4\)}
    {\(5\)}
    {\(6\)}
\end{problem}

\begin{problem}
在正方体 \(ABCD-A_1B_1C_1D_1\) 中,\(E\),\(F\) 分别为 \(AB\),\(BC\) 的中点,则(\quad)

\choices
    {\(\text{平面 }B_1EF\perp \text{平面 }BDD_1\)}
    {\(\text{平面 }B_1EF\perp \text{平面 }A_1BD\)}
    {\(\text{平面 }B_1EF // \text{平面 }A_1AC\)}
    {\(\text{平面 }B_1EF // \text{平面 }A_1C_1D\)}
\end{problem}

\begin{problem}
已知等比数列 \(\{a_n\}\) 的前 \(3\) 项和为 \(168\),\(a_2-a_5=42\),则 \(a_6=\)(\quad)

\choices
    {\(14\)}
    {\(12\)}
    {\(6\)}
    {\(3\)}
\end{problem}

\begin{problem}
已知球 \(O\) 的半径为 \(1\),四棱锥的顶点为 \(O\),底面的四个顶点均在球 \(O\) 的球面上,则当该四棱锥的体积最大时,其高为(\quad)

\choices
    {\(\frac{1}{3}\)}
    {\(\frac{1}{2}\)}
    {\(\frac{\sqrt{3}}{3}\)}
    {\(\frac{\sqrt{2}}{2}\)}
\end{problem}

\begin{problem}
某棋手与甲,乙,丙三位棋手各比赛一盘,各盘比赛结果相互独立. 已知该棋手与甲,乙,丙比赛获胜的概率分别为 \(p_1\),\(p_2\),\(p_3\),且 \(p_3>p_2>p_1>0\). 记该棋手连胜两盘的概率为 \(p\),则(\quad)

\choices
    {\(p\) 与该棋手和甲,乙,丙的比赛次序无关}
    {该棋手在第二盘与甲比赛,\(p\) 最大}
    {该棋手在第二盘与乙比赛,\(p\) 最大}
    {该棋手在第二盘与丙比赛,\(p\) 最大}
\end{problem}

\begin{problem}
双曲线 \(C\) 的两个焦点为 \(F_1\),\(F_2\),以 \(C\) 的实轴为直径的圆记为 \(D\),过 \(F_1\) 作 \(D\) 的切线与 \(C\) 交于 \(M\),\(N\) 两点,且 \(\cos \angle F_1NF_2=\frac{3}{5}\),则 \(C\) 的离心率为(\quad)

\choices
    {\(\frac{\sqrt{5}}{2}\)}
    {\(\frac{3}{2}\)}
    {\(\frac{\sqrt{13}}{2}\)}
    {\(\frac{\sqrt{17}}{2}\)}
\end{problem}

\begin{problem}
已知函数 \(f(x)\),\(g(x)\) 的定义域均为 \(\R\),且 \(f(x)+g(2-x)=5\),\(g(x)-f(x-4)=7\). 若 \(y=g(x)\) 的图像关于直线 \(x=2\) 对称,\(g(2)=4\),则 \(\sum_{k=1}^{22}f(k)=\)(\quad)

\choices
    {\(-21\)}
    {\(-22\)}
    {\(-23\)}
    {\(-24\)}
\end{problem}

\section{填空题}

\begin{problem}
从甲,乙等 \(5\) 名同学中随机选 \(3\) 名参加社区服务工作,则甲,乙都入选的概率为\fillinblank{}.
\end{problem}

\begin{problem}
过四点 \((0,0)\),\((4,0)\),\((-1,1)\),\((4,2)\) 中的三点的一个圆的方程为\fillinblank{}.
\end{problem}

\begin{problem}
记函数 \(f(x)=\cos \left( \omega x+\varphi \right)\)(\(\omega>0\),\(0<\varphi<\pi\))的最小正周期为 \(T\),若 \(f(T)=\frac{\sqrt{3}}{2}\),\(x=\frac{\pi}{9}\) 为 \(f(x)\) 的零点,则 \(\omega\) 的最小值为\fillinblank{}.
\end{problem}

\begin{problem}
已知 \(x=x_1\) 和 \(x=x_2\) 分别是函数 \(f(x)=2a^x-\e x^2\)(\(a>0\) 且 \(a\neq 1\))的极小值点和极大值点. 若 \(x_1<x_2\),则 \(a\) 的取值范围是\fillinblank{}.
\end{problem}

\section{解答题}

\begin{problem}
记 \(\triangle ABC\) 的内角 \(A\),\(B\),\(C\) 的对边分别为 \(a\),\(b\),\(c\),已知 \(\sin C\sin \left( A-B \right)=\sin B\sin \left( C-A \right)\).
\begin{enumerate}
\item 证明:\(2a^2=b^2+c^2\);
\item 若 \(a=5\),\(\cos A=\frac{25}{31}\),求 \(\triangle ABC\) 的周长.
\end{enumerate}
\end{problem}

\begin{problem}
如图,四面体 \(ABCD\) 中,\(AD\perp CD\),\(AD=CD\),\(\angle ADB=\angle BDC\),\(E\) 为 \(AC\) 的中点.

\begin{enumerate}
\item 证明:\(\text{平面 }BED\perp \text{平面 }ACD\);
\item 设 \(AB=BD=2\),\(\angle ACB=60^\circ\),点 \(F\) 在 \(BD\) 上. 当 \(\triangle AFC\) 面积最小时,求 \(CF\) 与平面 \(ABD\) 所成的角的正弦值.
\end{enumerate}

\bitmapfigure[max width=6.0cm]{img/2022/national_paper_b_science/q18_fig1.png}
\end{problem}

\begin{problem}
某地经过多年的环境治理,已将荒山改造成了绿水青山. 为估计一林区某种树木的总材积量,随机选取了 \(10\) 棵这种树木,测量每棵树的根部横截面积(单位:\( \mathrm{m}^2 \))和材积量(单位:\( \mathrm{m}^3 \)),得到如下数据:

\begin{center}
\renewcommand{\arraystretch}{1.3}
\begin{tabular}{|c|c|c|c|c|c|}
\hline
样本号 \(i\) & 1 & 2 & 3 & 4 & 5 \\
\hline
根部横截面积 \(x_i\) & 0.04 & 0.06 & 0.04 & 0.08 & 0.08 \\
\hline
材积量 \(y_i\) & 0.25 & 0.40 & 0.22 & 0.54 & 0.51 \\
\hline
\end{tabular}
\end{center}

\begin{center}
\renewcommand{\arraystretch}{1.3}
\begin{tabular}{|c|c|c|c|c|c|c|}
\hline
样本号 \(i\) & 6 & 7 & 8 & 9 & 10 & 总和 \\
\hline
根部横截面积 \(x_i\) & 0.05 & 0.05 & 0.07 & 0.07 & 0.06 & 0.6 \\
\hline
材积量 \(y_i\) & 0.34 & 0.36 & 0.46 & 0.42 & 0.40 & 3.9 \\
\hline
\end{tabular}
\end{center}

并计算得 \( \sum_{i=1}^{10}x_i^2=0.038 \),\( \sum_{i=1}^{10}y_i^2=1.6158 \),\( \sum_{i=1}^{10}x_iy_i=0.2474 \).

附:相关系数 \( r=\frac{\sum_{i=1}^{n}\left( x_i-\bar{x} \right)\left( y_i-\bar{y} \right)}{\sqrt{\sum_{i=1}^{n}\left( x_i-\bar{x} \right)^2\sum_{i=1}^{n}\left( y_i-\bar{y} \right)^2}} \),\( \sqrt{1.896}\approx1.377 \).

\begin{enumerate}
\item 估计该林区这种树木平均一棵的根部横截面积与平均一棵的材积量;
\item 求该林区这种树木的根部横截面积与材积量的样本相关系数(精确到 \(0.01\));
\item 现测量了该林区所有这种树木的根部横截面积,并得到所有这种树木的根部横截面积总和为 \(186\mathrm{m}^2\). 已知树木的材积量与其根部横截面积近似成正比. 利用以上数据给出该林区这种树木的总材积量的估计值.
\end{enumerate}
\end{problem}

\begin{problem}
已知椭圆 \(E\) 的中心为坐标原点,对称轴为 \(x\) 轴,\(y\) 轴,且过 \(A(0,-2)\),\(B\left( \frac{3}{2},-1 \right)\) 两点.
\begin{enumerate}
\item 求 \(E\) 的方程;
\item 设过点 \(P(1,-2)\) 的直线交 \(E\) 于 \(M\),\(N\) 两点,过 \(M\) 且平行于 \(x\) 轴的直线与线段 \(AB\) 交于点 \(T\),点 \(H\) 满足 \(MT=TH\). 证明:直线 \(HN\) 过定点.
\end{enumerate}
\end{problem}

\begin{problem}
已知函数 \(f(x)=\ln(1+x)+ax\e^{-x}\).
\begin{enumerate}
\item 当 \(a=1\) 时,求曲线 \(y=f(x)\) 在点 \(\left( 0,f(0) \right)\) 处的切线方程;
\item 若 \(f(x)\) 在区间 \((-1,0)\),\((0,+\infty)\) 各恰有一个零点,求 \(a\) 的取值范围.
\end{enumerate}
\end{problem}

\section{选考题:解答题}

\begin{problem}
在直角坐标系 \(xOy\) 中,曲线 \(C\) 的参数方程为 \(x=\sqrt{3}\cos 2t\),\(y=2\sin t\)(\(t\) 为参数),以坐标原点为极点,\(x\) 轴正半轴为极轴建立极坐标系,已知直线 \(l\) 的极坐标方程为 \( \rho \sin \left( \theta+\frac{\pi}{3} \right)+m=0 \).
\begin{enumerate}
\item 写出 \(l\) 的直角坐标方程;
\item 若 \(l\) 与 C 有公共点,求 \(m\) 的取值范围.
\end{enumerate}
\end{problem}

\begin{problem}
已知 \(a\),\(b\),\(c\) 都是正数,且 \(a^{\frac{3}{2}}+b^{\frac{3}{2}}+c^{\frac{3}{2}}=1\),证明:
\begin{enumerate}
\item \(abc\le \frac{1}{9}\);
\item \(\frac{a}{b+c}+\frac{b}{a+c}+\frac{c}{a+b}\le \frac{1}{2\sqrt{abc}}\).
\end{enumerate}
\end{problem}

